% 第 2 章重绘入口：课程案例现状与目标蓝图严格区分。
% 007 模型图共通样式；需在主文档已载入 TikZ、xcolor 后 input。
\tikzset{
  bsrActor/.style={draw=Navy,fill=white,line width=.8pt,rounded corners=1pt,align=center,minimum width=23mm,minimum height=8mm,font=\small},
  bsrWorker/.style={draw=Teal,fill=IceLight,line width=.8pt,rounded corners=1.2mm,align=center,minimum width=26mm,minimum height=8mm,font=\small},
  bsrSystem/.style={draw=Teal,fill=Glacier!45,line width=1pt,rounded corners=1.2mm,align=center,minimum width=28mm,minimum height=8mm,font=\small\bfseries},
  bsrEntity/.style={draw=Muted,fill=white,line width=.7pt,align=center,minimum width=26mm,minimum height=8mm,font=\small},
  bsrNote/.style={draw=Amber!85!Ink,fill=GoldPale,rounded corners=1mm,align=left,font=\scriptsize,text width=45mm,inner sep=3pt},
  bsrMessage/.style={-{Latex[length=2mm]},draw=Navy,line width=.75pt},
  bsrUnknown/.style={-{Latex[length=1.7mm]},draw=Amber!75!Ink,dashed,line width=.7pt},
  bsrReturn/.style={-{Latex[length=1.7mm]},draw=Muted,dashed,line width=.6pt},
  bsrState/.style={draw=Teal,fill=IceLight,rounded corners=2mm,line width=.8pt,align=center,font=\small,minimum height=8mm},
  bsrClass/.style={draw=Navy,fill=white,line width=.8pt,rectangle split,rectangle split parts=3,align=left,font=\scriptsize,inner sep=3pt},
  bsrER/.style={draw=Teal,fill=white,line width=.8pt,rectangle split,rectangle split parts=2,align=left,font=\scriptsize,inner sep=3pt}
}

% 业务序列图沿用课程课件的图例：执行者=小人，业务工人=圈中小人，
% 业务实体=空圆圈；工人与实体的右下斜线用于区分业务建模对象。
% 名称独立放在图标下方，以免首条消息压住图标或标题。
\newcommand{\BSRBizActor}[2]{%
  \begin{scope}[shift={(#1,0)}]
    \draw[Navy,line width=.85pt] (0,.72) circle[radius=1.2mm];
    \draw[Navy,line width=.85pt] (0,.56)--(0,.19) (-.22,.43)--(.22,.43)
      (0,.19)--(-.18,-.08) (0,.19)--(.18,-.08);
    \node[anchor=north,align=center,text width=2.9cm,font=\scriptsize] at (0,-.28) {#2};
  \end{scope}%
}
\newcommand{\BSRBizWorker}[2]{%
  \begin{scope}[shift={(#1,0)}]
    \draw[Teal,line width=.9pt,fill=IceLight] (0,.4) circle[radius=4.4mm];
    \draw[Teal,line width=.75pt] (0,.65) circle[radius=1mm];
    \draw[Teal,line width=.75pt] (0,.52)--(0,.27) (-.15,.41)--(.15,.41)
      (0,.27)--(-.13,.1) (0,.27)--(.13,.1);
    \draw[Teal,line width=.9pt] (.13,-.13)--(.37,.12);
    \node[anchor=north,align=center,text width=2.9cm,font=\scriptsize] at (0,-.28) {#2};
  \end{scope}%
}
\newcommand{\BSRBizEntity}[2]{%
  \begin{scope}[shift={(#1,0)}]
    \draw[Muted,line width=.9pt,fill=white] (0,.4) circle[radius=4.4mm];
    \draw[Muted,line width=.9pt] (.13,-.13)--(.37,.12);
    \node[anchor=north,align=center,text width=2.9cm,font=\scriptsize] at (0,-.28) {#2};
  \end{scope}%
}

% 四泳道业务序列图骨架：动作由第五参数给出；虚线表示待访谈/未证实。
\newcommand{\BSRBizSeq}[5]{%
\begin{tikzpicture}[x=1cm,y=.75cm]
  \BSRBizActor{0}{#1}
  \BSRBizWorker{4.3}{#2}
  \BSRBizEntity{8.6}{#3}
  \BSRBizEntity{12.9}{#4}
  \foreach \xx in {0,4.3,8.6,12.9} {\draw[Muted,densely dotted,line width=.45pt] (\xx,-1.28) -- (\xx,-7.75);}
  \begin{scope}[yshift=-1cm]#5\end{scope}
\end{tikzpicture}%
}

% 五泳道仅用于现状维修：分包商是业务工人，装备是独立业务实体。
\newcommand{\BSRBizSeqFive}[6]{%
\begin{tikzpicture}[x=1cm,y=.75cm]
  \BSRBizActor{0}{#1}
  \BSRBizWorker{3.225}{#2}
  \BSRBizEntity{6.45}{#3}
  \BSRBizWorker{9.675}{#4}
  \BSRBizEntity{12.9}{#5}
  \foreach \xx in {0,3.225,6.45,9.675,12.9} {\draw[Muted,densely dotted,line width=.45pt] (\xx,-1.28) -- (\xx,-7.75);}
  \begin{scope}[yshift=-1cm]#6\end{scope}
\end{tikzpicture}%
}

% 分模块三段式 UML 类框：id、x、y、名称、属性、操作。
\newcommand{\BSRClassCard}[6]{%
  \node[bsrClass,text width=3.45cm] (#1) at (#2,#3)
  {\centering\bfseries #4\nodepart{second}\raggedright #5\nodepart{third}\raggedright #6};%
}
\newcommand{\BSRERCard}[5]{%
  \node[bsrER,text width=3.65cm] (#1) at (#2,#3)
  {\centering\bfseries #4\nodepart{second}\raggedright #5};%
}

% 分析序列：执行者 + 边界类 + 控制类 + 两个实体类。
\newcommand{\BSRAnalysisSeq}[6]{%
\begin{tikzpicture}[x=1cm,y=.75cm]
  \node[bsrActor] at (0,0) {#1};
  \node[bsrWorker] at (3.5,0) {\texttt{<<boundary>>}\\#2};
  \node[bsrSystem] at (7,0) {\texttt{<<control>>}\\#3};
  \node[bsrEntity] at (10.5,0) {\texttt{<<entity>>}\\#4};
  \node[bsrEntity] at (14,0) {\texttt{<<entity>>}\\#5};
  \foreach \xx in {0,3.5,7,10.5,14} {\draw[Muted,densely dotted,line width=.45pt] (\xx,-.75) -- (\xx,-8.0);}
  \begin{scope}[yshift=-.55cm]#6\end{scope}
\end{tikzpicture}%
}

% 双执行者用例单列两个岗位，避免将顾客/员工或技师/经理混作同一责任人。
\newcommand{\BSRAnalysisSeqSix}[7]{%
\begin{tikzpicture}[x=1cm,y=.75cm]
  \node[bsrActor,minimum width=20mm,font=\scriptsize] at (0,0) {#1};
  \node[bsrActor,minimum width=20mm,font=\scriptsize] at (2.8,0) {#2};
  \node[bsrWorker,minimum width=22mm,font=\scriptsize] at (5.6,0) {\texttt{<<boundary>>}\\#3};
  \node[bsrSystem,minimum width=22mm,font=\scriptsize] at (8.4,0) {\texttt{<<control>>}\\#4};
  \node[bsrEntity,minimum width=22mm,font=\scriptsize] at (11.2,0) {\texttt{<<entity>>}\\#5};
  \node[bsrEntity,minimum width=22mm,font=\scriptsize] at (14,0) {\texttt{<<entity>>}\\#6};
  \foreach \xx in {0,2.8,5.6,8.4,11.2,14} {\draw[Muted,densely dotted,line width=.45pt] (\xx,-.75) -- (\xx,-9.8);}
  \begin{scope}[yshift=-.55cm]#7\end{scope}
\end{tikzpicture}%
}

\newcommand{\BSRStyleSample}{\BSRBizSeq{顾客\\业务执行者}{租赁专员\\业务工人}{BSR 系统\\业务实体}{物理单件\\业务实体}{%
  \draw[bsrMessage] (0,-1.1) -- node[above,font=\scriptsize]{选择租期与尺码} (8.6,-1.1);
  \draw[bsrMessage] (8.6,-2.1) -- node[above,font=\scriptsize]{分派备货} (4.3,-2.1);
  \draw[bsrMessage] (4.3,-3.1) -- node[above,font=\scriptsize]{现场核验} (12.9,-3.1);
  \draw[bsrReturn] (8.6,-4.1) -- node[above,font=\scriptsize]{准备进度} (0,-4.1);
  \node[bsrNote] at (6.45,-5.5) {图例：实线为请求/实物动作；虚线为返回或通知。日期与数量由系统核验，实物安全由员工确认。};
}}

% 研究对象为 BSR 组织；内部作业是业务用例的实现，不单列为顾客用例。
\newcommand{\BSRBusinessCases}{%
\begin{tikzpicture}[x=1cm,y=1cm]
  \draw[Teal,line width=.8pt,rounded corners=2mm,fill=IceLight] (2.1,-.5) rectangle (12.5,-5.2);
  \node[anchor=north west,font=\small\bfseries,text=Navy] at (2.4,-.7) {BSR 滑雪租赁业务边界};
  \begin{scope}[shift={(0,-1.8)}]\BSRBizActor{0}{顾客\\外部业务执行者}\end{scope}
  \node[draw=Teal,fill=white,ellipse,align=center,minimum width=6.4cm,minimum height=1cm] (rent) at (7.3,-1.9) {BU-01 租到合适、安全的装备};
  \node[draw=Teal,fill=white,ellipse,align=center,minimum width=6.4cm,minimum height=1cm] (settle) at (7.3,-3.5) {BU-02 归还装备并取得结算};
  \draw[Navy,line width=.7pt] (.3,-1.35) -- (rent.west);
  \draw[Navy,line width=.7pt] (.3,-1.35) -- (settle.west);
  \draw[Teal,line width=.7pt] (9.8,-2.22) -- (10.2,-1.95);
  \draw[Teal,line width=.7pt] (9.8,-3.82) -- (10.2,-3.55);
  \node[align=center,text width=9.4cm,font=\scriptsize,text=Muted] at (7.3,-4.65) {内部维修、备货与调拨支撑上述价值，属于业务用例的实现。};
  \node[draw=Muted,fill=white,rounded corners=1mm,align=left,text width=12cm,font=\scriptsize,inner sep=3mm] at (6.3,-6.1)
    {经营负责人：组织内部关键涉众，参与批准愿景与政策。\\
     度假村：组织外部关键涉众；现有证据未确立其直接业务交互，故不连接业务用例。};
\end{tikzpicture}%
}

% 仅表达案例可支持的事实片段；缺证据处使用说明卡，不补造业务顺序。
\newcommand{\BSRAsIsBooking}{\BSRBizDiagram{%
  \BSRBizWorker{0}{BSR 员工\\业务工人}
  \BSRBizEntity{12.9}{Excel 软件\\业务实体}
}{0,12.9}{%
  \draw[bsrMessage] (0,-1.8) -- node[above,font=\scriptsize]{管理客户、装备和合同信息（工作簿）} (12.9,-1.8);
}{-2.6}}

\newcommand{\BSRAsIsHandover}{\BSRBizEvidenceCard{B-ASIS-02 / EV-02：现状证据空缺}{%
案例要求逐件保存借出时间，但未提供现行取货、核身份、检查和交付的步骤与岗位分工。\\[1mm]
本卡为证据空缺说明，\textbf{不是业务序列图}；取得现场观察或访谈记录后再绘制现状协作。}}

\newcommand{\BSRAsIsReturn}{\BSRBizDiagram{%
  \BSRBizActor{0}{顾客\\业务执行者}
  \BSRBizWorker{12.9}{接收员工\\业务工人}
}{0,12.9}{%
  \draw[bsrMessage] (0,-1.8) -- node[above,font=\scriptsize]{受理异店归还（装备）} (12.9,-1.8);
}{-2.6}}

\newcommand{\BSRAsIsRepair}{\BSRBizSeqFive{门店顾客\\执行者}{BSR 员工\\工人}{纸质维修单\\业务实体}{分包商\\工人}{待修装备\\业务实体}{%
  \draw[bsrMessage] (3.225,-1.15) -- node[above,font=\tiny]{发现问题，填写纸单（EV-04）} (6.45,-1.15);
  \draw[bsrMessage] (3.225,-2.25) -- node[above,font=\tiny]{纸单附实物，移至维修架（EV-04）} (12.9,-2.25);
  \draw[bsrMessage] (9.675,-3.35) -- node[above,font=\tiny]{每周到访/带回工坊（EV-04）} (12.9,-3.35);
  \draw[bsrUnknown] (9.675,-4.45) -- node[above,font=\tiny]{结果/再上架确认（待访谈）} (3.225,-4.45);
  \node[bsrNote] at (6.45,-5.5) {旺季设备可能滞留一周以上。纸单流转可证实；修后安全验收责任未见案例描述。};
}}

\newcommand{\BSRAsIsTransfer}{\BSRBizSeq{顾客\\执行者}{门店员工\\工人}{实体装备\\业务实体}{门店/Excel\\业务实体}{%
  \draw[bsrUnknown] (0,-1.15) -- node[above,font=\scriptsize]{偶有异店归还（EV-05）} (4.3,-1.15);
  \draw[bsrUnknown] (4.3,-2.25) -- node[above,font=\scriptsize]{接收与记时（待访谈）} (8.6,-2.25);
  \draw[bsrUnknown] (4.3,-3.35) -- node[above,font=\scriptsize]{调拨批准/记账（待访谈）} (12.9,-3.35);
  \draw[bsrUnknown] (12.9,-4.45) -- node[above,font=\scriptsize]{运输和接收（待访谈）} (4.3,-4.45);
  \node[bsrNote] at (6.45,-5.5) {案例只说明归还地点偶有不同及永久归属可变；现行调拨流程未描述。};
}}

% 目标业务蓝图：只画人员与 BSR 软件间的职责协作，政策待 BSR 确认。
\newcommand{\BSRToBeBooking}{\BSRBizDiagram{%
  \BSRBizActor{0}{顾客\\业务执行者}
  \BSRBizEntity{6.45}{BSR 系统\\业务实体}
  \BSRBizWorker{12.9}{租赁专员\\业务工人}
}{0,6.45,12.9}{%
  \draw[bsrMessage] (0,-1.8) -- node[above,font=\scriptsize]{查询可租装备（门店、租期、尺码）} (6.45,-1.8);
  \draw[bsrReturn] (6.45,-2.8) -- node[above,font=\scriptsize]{可订量与 NOK 分项报价} (0,-2.8);
  \draw[bsrMessage] (0,-3.8) -- node[above,font=\scriptsize]{同意规则并申请暂留（预约明细）} (6.45,-3.8);
  \draw[bsrReturn] (6.45,-4.8) -- node[above,font=\scriptsize]{有限期暂留与当前合同版本} (0,-4.8);
  \draw[bsrMessage] (0,-5.8) -- node[above,font=\scriptsize]{本人确认当前合同版本} (6.45,-5.8);
  \draw[bsrMessage] (0,-6.8) -- node[above,font=\scriptsize]{完成本地模拟付款} (6.45,-6.8);
  \draw[bsrNotify] (6.45,-7.8) -- node[above,font=\scriptsize]{付款成功后分派具名备货任务} (12.9,-7.8);
}{-8.6}}

\newcommand{\BSRToBeHandover}{\BSRBizDiagram{%
  \BSRBizActor{0}{顾客\\业务执行者}
  \BSRBizWorker{6.45}{租赁专员\\业务工人}
  \BSRBizEntity{12.9}{BSR 系统\\业务实体}
}{0,6.45,12.9}{%
  \draw[bsrNotify] (12.9,-1.8) -- node[above,font=\scriptsize]{通知已付款的具名备货任务} (6.45,-1.8);
  \BSRBizSelf{6.45}{-2.9}{逐件备货并检查（装备）}
  \draw[bsrMessage] (6.45,-4) -- node[above,font=\scriptsize]{更新具名备货进度（检查事实）} (12.9,-4);
  \draw[bsrNotify] (12.9,-5) -- node[above,font=\scriptsize]{同步准备进度与当前处理人} (0,-5);
  \draw[bsrMessage] (0,-6) -- node[above,font=\scriptsize]{办理取货（取货码、身份）} (6.45,-6);
  \draw[bsrMessage] (6.45,-7) -- node[above,font=\scriptsize]{核验可交付资格（订单、合同）} (12.9,-7);
  \draw[bsrReturn] (12.9,-8) -- node[above,font=\scriptsize]{当前合同确认及付款核验结果} (6.45,-8);
  \draw[bsrMessage] (6.45,-9) -- node[above,font=\scriptsize]{验收并领取（合格装备）} (0,-9);
  \draw[bsrMessage] (6.45,-10) -- node[above,font=\scriptsize]{确认逐件交付（单件、处理人、时刻）} (12.9,-10);
}{-10.8}}

\newcommand{\BSRToBeReturn}{\BSRBizDiagram{%
  \BSRBizActor{0}{顾客\\业务执行者}
  \BSRBizWorker{6.45}{接收专员\\业务工人}
  \BSRBizEntity{12.9}{BSR 系统\\业务实体}
}{0,6.45,12.9}{%
  \draw[bsrMessage] (0,-1.8) -- node[above,font=\scriptsize]{受理原店或异店归还（部分实物）} (6.45,-1.8);
  \BSRBizSelf{6.45}{-2.9}{逐件检查（实物、库存号）}
  \draw[bsrMessage] (6.45,-4) -- node[above,font=\scriptsize]{登记归还（实还时刻、接收店、事实）} (12.9,-4);
  \draw[bsrReturn] (12.9,-5) -- node[above,font=\scriptsize]{本批受理结果与异常处置要求} (6.45,-5);
  \BSRBizSelf{6.45}{-6.1}{隔离疑损装备，交待复核}
  \draw[bsrNotify] (12.9,-7.2) -- node[above,font=\scriptsize]{通知归还结果及待结算事项} (0,-7.2);
}{-8}}

\newcommand{\BSRToBeRepair}{\BSRBizDiagram{%
  \BSRBizWorker{0}{接收专员\\业务工人}
  \BSRBizEntity{4.3}{BSR 系统\\业务实体}
  \BSRBizWorker{8.6}{装备检修技师\\业务工人}
  \BSRBizWorker{12.9}{门店经理\\业务工人}
}{0,4.3,8.6,12.9}{%
  \draw[bsrMessage] (0,-1.8) -- node[above,font=\scriptsize]{登记隔离故障（单件、事实）} (4.3,-1.8);
  \draw[bsrMessage] (12.9,-2.8) -- node[above,font=\scriptsize]{授权指派（技师、费率版本）} (4.3,-2.8);
  \draw[bsrNotify] (4.3,-3.8) -- node[above,font=\scriptsize]{通知初评与维修任务} (8.6,-3.8);
  \BSRBizSelf{8.6}{-4.9}{诊断拍照（装备受损部位）}
  \draw[bsrMessage] (8.6,-6) -- node[above,font=\scriptsize]{24h 内提交初评（级别、照片、依据）} (4.3,-6);
  \BSRBizSelf{4.3}{-7.1}{发布顾客可见初评与处理人}
  \BSRBizSelf{8.6}{-8.2}{按授权维修并安全自检（装备）}
  \draw[bsrMessage] (8.6,-9.3) -- node[above,font=\scriptsize]{申请验收（工时、材料、安全项）} (4.3,-9.3);
  \draw[bsrNotify] (4.3,-10.3) -- node[above,font=\scriptsize]{通知待独立安全验收} (12.9,-10.3);
  \BSRBizSelfLeft{12.9}{-11.4}{实物安全验收（维修后装备）}
  \draw[bsrMessage] (12.9,-12.5) -- node[above,font=\scriptsize]{记录复租、返修或报废结论} (4.3,-12.5);
}{-13.2}}

\newcommand{\BSRToBeTransfer}{\BSRBizDiagram{%
  \BSRBizWorker{0}{调出店专员\\业务工人}
  \BSRBizWorker{4.3}{门店经理\\业务工人}
  \BSRBizEntity{8.6}{BSR 系统\\业务实体}
  \BSRBizWorker{12.9}{接收店专员\\业务工人}
}{0,4.3,8.6,12.9}{%
  \draw[bsrMessage] (0,-1.8) -- node[above,font=\scriptsize]{提出永久调拨（单件、目标店）} (8.6,-1.8);
  \draw[bsrMessage] (4.3,-2.8) -- node[above,font=\scriptsize]{批准调拨（可调单件）} (8.6,-2.8);
  \draw[bsrNotify] (8.6,-3.8) -- node[above,font=\scriptsize]{通知调出并标记在途} (0,-3.8);
  \BSRBizSelf{0}{-4.9}{组织运输（实物装备）}
  \draw[bsrMessage] (0,-6) -- node[above,font=\scriptsize]{接收交接（实物、调拨凭据）} (12.9,-6);
  \BSRBizSelfLeft{12.9}{-7.1}{逐件验收（到店实物）}
  \draw[bsrMessage] (12.9,-8.2) -- node[above,font=\scriptsize]{确认接收，更新永久归属} (8.6,-8.2);
}{-9}}


% 007 业务模型文本：正文第 2 章输入。
\subsection{2.1 愿景}
\subsubsection{2.1.1 愿景与证据边界}
\BSRKeyPoint{组织愿景}{待 BSR 确认}{在第二门店及未来更多门店共用可信的客户、合同和单件库存记录，使顾客按预约时段快速取得合适、安全的装备，并让跨店归还和维修可追踪。}

\noindent 组织的关键涉众是经营负责人 Ole Fjordal 与度假村管理方。案例能证明第二租赁点扩张压力、缓解主店拥堵的意图和 Excel 单人操作局限；尚无 BSR 提供的排队、库存差错或维修周转基线。下表是建议指标，不能写成已完成绩效。

\begin{BSRTable}
\rowcolor{TableHead}\color{white}\bfseries 指标及口径 & \color{white}\bfseries 建议目标与状态 \\
预约可履约率：可用单件覆盖已确认预约的比例 & 建议 $\geq 98\%$；基线与目标待 BSR 确认。\\
到店取货时长：核身份至完成逐件交付的中位数 & 建议 $\leq 10$ 分钟；须现场测量。\\
库存账实一致率：盘点一致实物数/抽盘实物数 & 建议 $\geq 99\%$；须跨店抽盘验证。\\
疑似损坏初评及时率：发现后 24 小时内反馈的比例 & 建议 $\geq 95\%$；初评不是维修完成。\\
\end{BSRTable}

\subsubsection{2.1.2 执行者、工人与业务实体}
\begin{BSRTable}
\rowcolor{TableHead}\color{white}\bfseries 身份 & \color{white}\bfseries 价值、职责与证据边界 \\
业务执行者：顾客 & 取得合适装备、透明租金与押金、按约取还及可核对的异常费用。线上预约属目标蓝图。\\
内部关键涉众：经营负责人 Ole & 代表 BSR 经营目标，批准愿景、费率与跨店政策；其管理岗位在组织内部，不列为外部业务执行者。\\
外部关键涉众：度假村管理方 & 关注主店拥堵和第二租赁点扩张；现有证据未说明其向 BSR 请求何种具体业务服务，暂不连接业务用例。\\
业务工人：租赁专员 & 案例支持员工参与合同及故障登记；目标中负责备货、核身份/码、交付与逐件归还。\\
业务工人：门店经理 & 目标中负责排班、调拨、受损复核和安全验收；现行分工待访谈。\\
业务工人：装备检修技师/分包商 & 案例明确分包商定期到访并按纸单维修；在所建模的维修责任范围内作为业务工人，目标中按授权工单记录证据与工时。\\
业务实体：Excel 软件、目标 BSR 系统 & 表示能承担信息处理职责的非人系统；现状仅有 Excel 用途证据，BSR 系统属于改进蓝图。\\
业务资料与实物：客户记录、合同、装备、纸质维修单、货架 & 是处理对象或消息参数；不赋予纸张、装备或货架业务处理责任，不将它们放在业务序列生命线上。\\
\end{BSRTable}

\subsection{2.2 业务用例图}
\subsubsection{2.2.1 对外业务用例与价值}
\begin{BSRTable}
\rowcolor{TableHead}\color{white}\bfseries 业务用例 & \color{white}\bfseries 对执行者的可观察结果 \\
BU-01 租到合适、安全的装备 & 顾客在约定门店和时段取得适合尺码与水平的安全装备，知悉租金、押金及租赁责任。预约、报价、合同确认、付款与备货交付共同实现该价值。\\
BU-02 归还装备并取得结算 & 顾客得到逐件归还确认与可核对的租金、超时费、损坏费及押金结果；异常或争议有明确处理状态。异店归还属于该服务的约定方式。\\
\end{BSRTable}

\begin{figure}[htbp]\centering\resizebox{.94\linewidth}{!}{\BSRBusinessCases}
\caption{BU-01～02 BSR 对顾客提供的完整业务价值；内部维修与调拨属于实现支撑}\label{fig:b-cases}\end{figure}\FloatBarrier

\noindent B-CASE-01～05 保留为旧版实现片段索引，便于核对既有追踪材料，\textbf{不再表示五个独立的对外业务用例}。
\begin{BSRTable}
\rowcolor{TableHead}\color{white}\bfseries 原实现片段索引 & \color{white}\bfseries 现行业务价值映射 \\
B-CASE-01 预约与报价 & 实现 BU-01；对应 B-ASIS-01 / B-TOBE-01。\\
B-CASE-02 备货与交付 & 实现 BU-01；对应 B-ASIS-02 / B-TOBE-02。\\
B-CASE-03 接收归还 & 实现 BU-02；对应 B-ASIS-03 / B-TOBE-03。\\
B-CASE-04 维护装备 & 内部检修支撑 BU-01，受损结算证据支撑 BU-02；对应 B-ASIS-04 / B-TOBE-04。\\
B-CASE-05 跨店作业 & 异店归还实现 BU-02；永久调拨为保障 BU-01 的内部库存作业，对应 B-ASIS-05 / B-TOBE-05。\\
\end{BSRTable}

\subsubsection{2.2.2 图示规范与证据方法}
\noindent 本章依《软件方法：业务建模和需求（上，第 2 版）》印刷页 68～78 区分组织外的业务执行者、组织内承担职责的业务工人，以及外部执行者期望取得的完整价值。组织边界按责任划分，不能只凭工作地点、雇佣或外包形式判断；经营负责人和度假村的愿景利益另列为涉众。

\noindent 依该书印刷页 101～105，业务序列中 A 指向 B 的消息表示要求 B 承担职责；生命线的粒度是人或非人系统。纸合同、维修单和装备没有处理能力，作为消息参数表达；检查、搬运和实物验收由人员自调用表达。图标保留课程中的小人、圈中小人和空圆圈，分别表示业务执行者、业务工人与非人业务实体。

\noindent 实线实箭头表示职责请求，实线开箭头表示主动通知；水平虚线箭头只表示返回消息，垂直虚线表示生命线。分析图的暖色实线表示异常分支，\textbf{线型不表示证据可信度}。现状仅绘案例已支持的事实片段，证据不足时给出说明卡并在证据表列出待核问题；不能用虚线拼成一套未经观察的现行流程。改进图均为待 BSR 确认的目标协作，系统自动核验和提醒不替代人员的实物责任。


\subsection{2.3 现状业务序列图}
\noindent 下列图的实线只承载案例明确的动作，暖色虚线表示尚无访谈/现场观察支持的候选步骤。现状图是要向 BSR 验证的工作底稿，不能反推为已发生的线上预约或自动退款。
\BSRNewFigure{\BSRAsIsBooking}{B-ASIS-01 现场开单现状：仅 Excel 工作簿用途有 EV-01 证据，其余待访谈}{fig:b-asis-01}
\BSRNewFigure{\BSRAsIsHandover}{B-ASIS-02 装备交付现状：顺序与职责待访谈（EV-02）}{fig:b-asis-02}
\BSRNewFigure{\BSRAsIsReturn}{B-ASIS-03 逐件与异店归还现状：异店现象有 EV-03 证据，处理步骤待访谈}{fig:b-asis-03}
\BSRNewFigure{\BSRAsIsRepair}{B-ASIS-04 装备维修现状：纸单和定期分包商有 EV-04 证据}{fig:b-asis-04}
\BSRNewFigure{\BSRAsIsTransfer}{B-ASIS-05 跨店转移现状：异店归还/归属可变有 EV-05 证据，调拨手续待访谈}{fig:b-asis-05}

\subsection{2.4 改进业务序列图}
\noindent 改进图是供 BSR 审阅的业务蓝图。系统把纸面/口头传递变为共享信息流，自动核验日期段库存、合同/政策版本、逐件计费和责任提醒；拣货、交接、搬运及安全验收保留人员责任。\textbf{蓝图未获 BSR 签字，不代表现行制度。}
\BSRNewFigure{\BSRToBeBooking}{B-TOBE-01 日期段预约与透明报价：目标蓝图，政策待 BSR 确认}{fig:b-tobe-01}
\BSRNewFigure{\BSRToBeHandover}{B-TOBE-02 具名备货、核码与逐件交付：目标蓝图}{fig:b-tobe-02}
\BSRNewFigure{\BSRToBeReturn}{B-TOBE-03 逐件异店归还与异常分流：目标蓝图}{fig:b-tobe-03}
\BSRNewFigure{\BSRToBeRepair}{B-TOBE-04 受损初评、维修与经理安全验收：目标蓝图}{fig:b-tobe-04}
\BSRNewFigure{\BSRToBeTransfer}{B-TOBE-05 跨店库存转移与接收：目标蓝图}{fig:b-tobe-05}

\subsubsection{2.4.1 蓝图问题与自动化边界}
\begin{BSRTable}
\rowcolor{TableHead}\color{white}\bfseries 原有问题/来源 & \color{white}\bfseries 改进及人工作业 \\
Excel 工作簿单人可靠操作（EV-01） & 顾客、两店员工读取同一服务端预约与库存事实；服务端按租期二次核量，员工仍需现场确认实物。\\
纸质维修单可能滞留（EV-04） & 电子工单与指派提醒缩短信息等待；技师提交后仍由经理实物验收。\\
异店归还与归属变更（EV-03/05） & 分开保存当前所在店与永久归属店；跨店运输和接收不能由系统虚拟完成。\\
押金、损坏和超时规则（设计需求） & 客户下单前知悉 NOK 分项与版本；受损费用需技师证据和经理复核，争议时暂停结清。\\
\end{BSRTable}
\subsubsection{2.4.2 网页与小程序共用的目标信息流}
\noindent 顾客无论从网页还是微信小程序进入，均向同一 BSR 服务端查询日期段库存、维护本人购物车并提交报价与预约；小程序首次使用须经过微信临时身份交换和现有顾客账号的显式验证绑定。服务端以同一顾客 ID、同一门店单件和同一订单进行核价与履约。租赁专员在网页端具名备货、验货交付和接收，技师/经理在网页端完成鉴定与复核；顾客两个客户端按本人权限看同一进度。正式微信扫码、支付与组织政策仍待外部条件和 BSR 确认。本段是改进蓝图，不是案例中的现状事实。
